\section{Experimental Setup}
\label{sec:experiments}

We design four interconnected experimental studies, each testing a distinct claim in our causal chain: orbits are large $\to$ NFT sees them $\to$ canonicalization concentrates structure $\to$ but the standard algorithm fails. The experiments are ordered to build evidence for this chain before revealing where it breaks.

\subsection{Research Questions}

\textbf{RQ1 (Orbit Existence, H-E1):} Do scaling and sign-flip symmetry orbits have non-negligible geometric diameter in the Schürholt MNIST MLP zoo --- and if so, how large are they?

\textbf{RQ2 (Encoder Invariance, H-M1):} Is NFT, trained on raw Schürholt zoo weights, naturally invariant to scaling and sign-flip orbits? Does it embed functionally identical networks similarly?

\textbf{RQ3 (Geometric Concentration, H-M2):} Does scaling canonicalization concentrate geometric structure in the weight space --- specifically, does it improve PCA explained variance and linear-probe property prediction?

\textbf{RQ4 (Uniqueness Audit, H-C1):} Is sign-flip canonicalization via the majority-sign algorithm well-defined for the Schürholt MNIST MLP architecture?

An implicit RQ5 (H-M3) evaluates end-to-end property prediction improvement under canonicalization, providing directional evidence and scale requirements.

\subsection{Dataset}

\textbf{Schürholt MNIST Model Zoo} \citep{Schurholt2022ModelZoos}. A collection of 2-layer MLPs (architecture: $784 \to 64 \to 10$, ReLU activations) trained on MNIST with varying learning rates (log-uniform in $[10^{-5}, 10^{-1}]$), weight decay values, and random initialization seeds.

\begin{table}[h]
\centering
\caption{Ground-truth property labels in the Schürholt MNIST zoo.}
\label{tab:properties}
\begin{tabular}{lll}
\toprule
Property & Description & Range (approx.) \\
\midrule
Test accuracy & MNIST test set accuracy & 0.60 -- 0.98 \\
Generalization gap & Train $-$ test accuracy & 0.00 -- 0.35 \\
Learning rate & Training LR (log-scale) & $10^{-5}$ -- $10^{-1}$ \\
\bottomrule
\end{tabular}
\end{table}

\textbf{Available subset.} We use a locally archived subset of $N=500$ models (full Schürholt zoo $\approx 50{,}000$ models; HuggingFace dataset unavailable at runtime). For orbit characterization (RQ1, RQ2), this subset is sufficient. For property prediction (RQ3, RQ5), $N=500$ is severely underpowered, as we quantify explicitly in Section~\ref{sec:results}.

\textbf{Weight representation.} Each model's weights are flattened to a vector of dimension $D=51{,}850$ ($784 \times 64 + 64 \times 10$). For oracle orbit construction, we retain the structured weight dictionary ($W_1, W_2$) to apply layer-specific transforms.

\subsection{Baselines}

\textbf{Raw weights $\to$ NFT (Condition A).} Standard NFT encoding without preprocessing. At $N=500$, NFT achieves Spearman $\rho \approx 0.11$ --- substantially below the layer statistics baseline ($\rho \approx 0.9$). This gap primarily reflects underpowered training ($N=400$ samples) combined with NFT devoting representational capacity to symmetry-induced weight variation rather than functional properties.

\textbf{Random normalization $\to$ NFT (Condition E).} Each hidden neuron's incoming weights are scaled by a random factor drawn from $\mathcal{N}(1, 0.1)$, preserving sign structure. Condition E is a \textit{null control}: if Condition D (canonical) outperforms E, the improvement is symmetry-specific; if $\mathrm{E} \geq \mathrm{D}$, the effect is a normalization artifact.

\textbf{Layer statistics} \citep{Unterthiner2020PredictingNN}. Mean, variance, and higher moments of weight distributions computed per layer. Achieves Spearman $\rho \approx 0.9$ on simple zoos. Included to contextualize the gap between NFT baseline and practical performance ceiling.

\subsection{Evaluation Metrics}

\textbf{Spearman $\rho$.} Rank correlation between predicted and ground-truth property labels on the 50-model held-out test split.

\textbf{Orbit invariance gap.} $\overline{\text{within-orbit cosine sim.}} - \overline{\text{cross-orbit same-property cosine sim.}}$ A positive gap indicates the encoder treats symmetry-related models as \textit{less} similar than property-matched unrelated models.

\textbf{Explained Variance Ratio (EVR).} Fraction of total weight-matrix variance captured by the top-$k$ principal components.

\textbf{Linear regression $R^2$.} Variance in property labels explained by the top-$k$ principal components.

\textbf{Fraction unique (H-C1).} Fraction of zoo models for which the majority-sign algorithm produces a well-defined unique canonical form.

\textbf{Statistical reporting.} All comparisons reported with bootstrap 95\% CIs ($n_{\mathrm{boot}}=1{,}000$). At $n=50$ test samples, Spearman $\rho$ CIs have width $\approx 0.6$.

\subsection{Implementation Details}

\textbf{NFT architecture.} $d_{\mathrm{model}}=256$, 4 self-attention layers, 8 heads, CLS-token pooling. Weights tokenized row-by-row. CLS-token pooling is critical: mean pooling collapses the embedding and destroys model discrimination (verified empirically). NFT has approximately 3.4M parameters.

\textbf{Training.} Adam optimizer, lr$=10^{-3}$, weight\_decay$=10^{-4}$. Maximum 50 epochs with early stopping on validation Spearman $\rho$ (patience$=10$).

\textbf{Data split.} 80/10/10 fixed split: 400 training, 50 validation, 50 test models. Fixed across conditions for fair comparison.

\textbf{Oracle orbit construction.} For each base model, one oracle orbit member constructed by applying the symmetry transform: $\alpha_i \sim \log\text{-Uniform}(0.1, 10.0)$ for scaling, $s_i \sim \mathrm{Bernoulli}(0.5) \to \{-1, +1\}$ for sign-flip. Seed fixed per experiment.

\textbf{Bootstrap CI.} \texttt{scipy.stats.bootstrap} with $n_{\mathrm{boot}}=1{,}000$, seed$=42$, 95\% confidence level.

\textbf{Hardware.} All experiments run on CPU. NFT training takes approximately 5--15 minutes per condition for $N=400$ training samples.
